% ====================================================================
% ФАЙЛ: tasks/02_mkt/mkt_fundamental_equations_bank_part1.tex
% ТЕМА: ОСНОВНОЕ УРАВНЕНИЕ МКТ. АБСОЛЮТНАЯ ТЕМПЕРАТУРА. КИНЕТИЧЕСКАЯ ЭНЕРГИЯ
% ПАЧКА 1: ВАРИАНТЫ 1–05
% ====================================================================

% === ЗАДАЧА 1 ===
% Тег: #МКТ #КИМ_09 #ВАРИАНТ_01
\begin{taskbasic}[code={\label{task:mkt_v1_n9}}]
\textbf{Задача 1.} Один моль разреженного аргона участвует в процессе 1–2–3, представленном на графике зависимости давления $p$ от средней кинетической энергии $\overline{E}_{\text{к}}$ теплового движения молекул.

\nopagebreak\smallskip
\begin{center}
\begin{tikzpicture}[scale=1.2, every node/.style={font=\footnotesize}]
    \draw[xstep=1, ystep=1, gray!30, thin] (0,0) grid (4,4);
    \draw[-{Stealth[scale=0.8]}, black, thick] (0,0) -- (4.5,0) node[right] {$\overline{E}_{\text{к}}$};
    \draw[-{Stealth[scale=0.8]}, black, thick] (0,0) -- (0,4.5) node[above] {$p$};
    
    \coordinate (1) at (1,1);
    \coordinate (2) at (3,2);
    \coordinate (3) at (3,4);
    
    \draw[dashed, thin, gray] (0,0) -- (1);
    \draw[ultra thick, brandPrimary, ->, >=stealth] (1) -- (2);
    \draw[ultra thick, brandPrimary, ->, >=stealth] (2) -- (3);
    
    \fill[black] (1) circle (1.5pt) node[below right] {1};
    \fill[black] (2) circle (1.5pt) node[right=2pt, yshift=-2pt] {2};
    \fill[black] (3) circle (1.5pt) node[above right] {3};
    \node[below left] at (0,0) {0};
\end{tikzpicture}
\end{center}

Из приведённого ниже списка выберите все верные утверждения, характеризующие процессы на рисунке.
1) В процессе 1–2 внутренняя энергия аргона уменьшается.
2) В процессе 2–3 аргон изотермически сжимают.
3) В процессе 1–2 концентрация аргона уменьшается.
4) В процессе 2–3 температура аргона растёт.
5) В процессе 2–3 аргон отдаёт окружающей среде положительное количество теплоты.

\kim{9}
\end{taskbasic}
\answer{25}

% === ЗАДАЧА 2 ===
% Тег: #МКТ #КИМ_09 #ВАРИАНТ_02
\begin{taskbasic}[code={\label{task:mkt_v2_n9}}]
\textbf{Задача 2.} Один моль разреженного аргона участвует в процессе 1–2–3, представленном на графике зависимости давления $p$ от средней кинетической энергии $\overline{E}_{\text{к}}$ теплового движения молекул.

\nopagebreak\smallskip
\begin{center}
\begin{tikzpicture}[scale=1.2, every node/.style={font=\footnotesize}]
    \draw[xstep=1, ystep=1, gray!30, thin] (0,0) grid (4,4);
    \draw[-{Stealth[scale=0.8]}, black, thick] (0,0) -- (4.5,0) node[right] {$\overline{E}_{\text{к}}$};
    \draw[-{Stealth[scale=0.8]}, black, thick] (0,0) -- (0,4.5) node[above] {$p$};
    
    \coordinate (1) at (1,1);
    \coordinate (2) at (3,2);
    \coordinate (3) at (3,4);
    
    \draw[dashed, thin, gray] (0,0) -- (1);
    \draw[ultra thick, brandPrimary, ->, >=stealth] (1) -- (2);
    \draw[ultra thick, brandPrimary, ->, >=stealth] (2) -- (3);
    
    \fill[black] (1) circle (1.5pt) node[below right] {1};
    \fill[black] (2) circle (1.5pt) node[right=2pt, yshift=-2pt] {2};
    \fill[black] (3) circle (1.5pt) node[above right] {3};
    \node[below left] at (0,0) {0};
\end{tikzpicture}
\end{center}

Из приведённого ниже списка выберите все верные утверждения, характеризующие процессы на рисунке.
1) В процессе 1–2 аргон получает положительное количество теплоты.
2) В процессе 2–3 аргон изотермически расширяется.
3) В процессе 1–2 концентрация аргона остаётся неизменной.
4) В процессе 2–3 температура аргона остаётся неизменной.
5) В процессе 2–3 внутренняя энергия аргона увеличивается.

\kim{9}
\end{taskbasic}
\answer{134}

% === ЗАДАЧА 3 ===
% Тег: #МКТ #КИМ_07 #ВАРИАНТ_03
\begin{taskbasic}[code={\label{task:mkt_v3_n7}}]
\textbf{Задача 3.} При уменьшении абсолютной температуры на $300$~К средняя кинетическая энергия поступательного теплового движения молекул неона уменьшилась в $1{,}5$ раза. Какова начальная температура газа?

\kim{7}
\end{taskbasic}
\answer{900}

% === ЗАДАЧА 4 ===
% Тег: #МКТ #КИМ_07 #ВАРИАНТ_04
\begin{taskbasic}[code={\label{task:mkt_v4_n7}}]
\textbf{Задача 4.} При уменьшении абсолютной температуры на $300$~К средняя кинетическая энергия поступательного теплового движения молекул аргона уменьшилась в $1{,}5$ раза. Какова конечная температура газа?

\kim{7}
\end{taskbasic}
\answer{600}

% === ЗАДАЧА 5 ===
% Тег: #МКТ #КИМ_07 #ВАРИАНТ_05
\begin{taskbasic}[code={\label{task:mkt_v5_n7}}]
\textbf{Задача 5.} В сосуде под поршнем находится некоторое постоянное количество разреженного неона. Во сколько раз увеличится среднеквадратическая скорость теплового движения молекул неона, если он перейдёт из состояния 1 в состояние 2 (см. рисунок)?

\nopagebreak\smallskip
\begin{center}
\begin{tikzpicture}[scale=1.0, every node/.style={font=\footnotesize}]
    \draw[xstep=1, ystep=1, gray!30, thin] (0,0) grid (5,5);
    \draw[-{Stealth[scale=0.8]}, black, thick] (0,0) -- (5.5,0) node[right] {$V$};
    \draw[-{Stealth[scale=0.8]}, black, thick] (0,0) -- (0,5.5) node[above] {$p$};
    
    \coordinate (1) at (1,3);
    \coordinate (2) at (3,4);
    
    \draw[ultra thick, brandPrimary, ->, >=stealth] (1) -- (2);
    
    \fill[black] (1) circle (1.5pt) node[above left] {1};
    \fill[black] (2) circle (1.5pt) node[above right] {2};
    
    \node[left] at (0,3) {$3p_0$};
    \node[left] at (0,4) {$4p_0$};
    \node[below] at (1,0) {$V_{0}$};
    \node[below] at (3,0) {$3V_{0}$};
    \node[below left] at (0,0) {0};
\end{tikzpicture}
\end{center}

\kim{7}
\end{taskbasic}
\answer{2}

% ====================================================================
% ФАЙЛ: tasks/02_mkt/mkt_fundamental_equations_bank_part2.tex
% ТЕМА: ОСНОВНОЕ УРАВНЕНИЕ МКТ. АБСОЛЮТНАЯ ТЕМПЕРАТУРА. КИНЕТИЧЕСКАЯ ЭНЕРГИЯ
% ПАЧКА 2: ВАРИАНТЫ 06–10
% ====================================================================

% === ЗАДАЧА 1 ===
% Тег: #МКТ #КИМ_07 #ВАРИАНТ_06
\begin{taskbasic}[code={\label{task:mkt_v6_n7}}]
\textbf{Задача 1.} В сосуде под поршнем находится некоторое постоянное количество разреженного аргона. Во сколько раз увеличится средняя кинетическая энергия теплового движения молекул аргона, если он перейдёт из состояния 1 в состояние 2 (см. рисунок)?

\nopagebreak\smallskip
\begin{center}
\begin{tikzpicture}[scale=1.0, every node/.style={font=\footnotesize}]
    \draw[xstep=1, ystep=1, gray!30, thin] (0,0) grid (5,5);
    \draw[-{Stealth[scale=0.8]}, black, thick] (0,0) -- (5.5,0) node[right] {$V$};
    \draw[-{Stealth[scale=0.8]}, black, thick] (0,0) -- (0,5.5) node[above] {$p$};
    
    \coordinate (1) at (1,2);
    \coordinate (2) at (3,4);
    
    \draw[ultra thick, brandPrimary, ->, >=stealth] (1) -- (2);
    
    \fill[black] (1) circle (1.5pt) node[above left] {1};
    \fill[black] (2) circle (1.5pt) node[above right] {2};
    
    \node[left] at (0,2) {$2p_0$};
    \node[left] at (0,4) {$4p_0$};
    \node[below] at (1,0) {$V_{0}$};
    \node[below] at (3,0) {$3V_{0}$};
    \node[below left] at (0,0) {0};
\end{tikzpicture}
\end{center}

\kim{7}
\end{taskbasic}
\answer{6}

% === ЗАДАЧА 2 ===
% Тег: #МКТ #КИМ_07 #ВАРИАНТ_09
\begin{taskbasic}[code={\label{task:mkt_v9_n7}}]
\textbf{Задача 2.} В сосуде под поршнем находится некоторое постоянное количество идеального газа. Во сколько раз увеличится средняя кинетическая энергия поступательного теплового движения молекул газа, если он перейдёт из состояния 2 в состояние 1 (см. рисунок)?

\nopagebreak\smallskip
\begin{center}
\begin{tikzpicture}[scale=1.0, every node/.style={font=\footnotesize}]
    \draw[xstep=1, ystep=1, gray!30, thin] (0,0) grid (5,5);
    \draw[-{Stealth[scale=0.8]}, black, thick] (0,0) -- (5.5,0) node[right] {$V$};
    \draw[-{Stealth[scale=0.8]}, black, thick] (0,0) -- (0,5.5) node[above] {$p$};
    
    \coordinate (1) at (4,2);
    \coordinate (2) at (2,4);
    
    \draw[ultra thick, brandPrimary, ->, >=stealth] (2) -- (1);
    
    \fill[black] (1) circle (1.5pt) node[above right] {1};
    \fill[black] (2) circle (1.5pt) node[above left] {2};
    
    \node[left] at (0,2) {$2p_0$};
    \node[left] at (0,4) {$4p_0$};
    \node[below] at (2,0) {$2V_{0}$};
    \node[below] at (4,0) {$4V_{0}$};
    \node[below left] at (0,0) {0};
\end{tikzpicture}
\end{center}

\kim{7}
\end{taskbasic}
\answer{1}

% === ЗАДАЧА 3 ===
% Тег: #МКТ #КИМ_07 #ВАРИАНТ_10
\begin{taskbasic}[code={\label{task:mkt_v10_n7}}]
\textbf{Задача 3.} В сосуде под поршнем находится некоторое постоянное количество идеального газа. Во сколько раз уменьшится температура газа, если он перейдёт из состояния 2 в состояние 1 (см. рисунок)?

\nopagebreak\smallskip
\begin{center}
\begin{tikzpicture}[scale=1.0, every node/.style={font=\footnotesize}]
    \draw[xstep=1, ystep=1, gray!30, thin] (0,0) grid (5,5);
    \draw[-{Stealth[scale=0.8]}, black, thick] (0,0) -- (5.5,0) node[right] {$V$};
    \draw[-{Stealth[scale=0.8]}, black, thick] (0,0) -- (0,5.5) node[above] {$p$};
    
    \coordinate (1) at (2,4);
    \coordinate (2) at (4,5); % В соответствии со значениями 8p0 и 4V0
    
    \draw[ultra thick, brandPrimary, ->, >=stealth] (2) -- (1);
    
    \fill[black] (1) circle (1.5pt) node[above left] {1};
    \fill[black] (2) circle (1.5pt) node[above right] {2};
    
    \node[left] at (0,4) {$4p_0$};
    \node[left] at (0,5) {$8p_0$};
    \node[below] at (2,0) {$2V_{0}$};
    \node[below] at (4,0) {$4V_{0}$};
    \node[below left] at (0,0) {0};
\end{tikzpicture}
\end{center}

\kim{7}
\end{taskbasic}
\answer{4}

% ====================================================================
% ФАЙЛ: tasks/02_mkt/mkt_fundamental_equations_bank_part3.tex
% ТЕМА: ОСНОВНОЕ УРАВНЕНИЕ МКТ. АБСОЛЮТНАЯ ТЕМПЕРАТУРА. КИНЕТИЧЕСКАЯ ЭНЕРГИЯ
% ПАЧКА 3: ВАРИАНТЫ 11–16
% ====================================================================

% === ЗАДАЧА 1 ===
% Тег: #МКТ #КИМ_07 #ВАРИАНТ_13
\begin{taskbasic}[code={\label{task:mkt_v13_n7}}]
\textbf{Задача 1.} В результате уменьшения абсолютной температуры средняя кинетическая энергия поступательного теплового движения молекул разреженного гелия уменьшилась в $4$ раза. Во сколько раз уменьшилась при этом абсолютная температура гелия?

\kim{7}
\end{taskbasic}
\answer{4}

% === ЗАДАЧА 2 ===
% Тег: #МКТ #КИМ_07 #ВАРИАНТ_14
\begin{taskbasic}[code={\label{task:mkt_v14_n7}}]
\textbf{Задача 2.} При неизменной концентрации молекул разреженного неонового газа средняя кинетическая энергия поступательного теплового движения его молекул увеличилась в $2$ раза. Во сколько раз увеличилось при этом давление газа в сосуде?

\kim{7}
\end{taskbasic}
\answer{2}

% ====================================================================
% ФАЙЛ: tasks/02_mkt/mkt_fundamental_equations_bank_part4.tex
% ТЕМА: ОСНОВНОЕ УРАВНЕНИЕ МКТ. АБСОЛЮТНАЯ ТЕМПЕРАТУРА. КИНЕТИЧЕСКАЯ ЭНЕРГИЯ
% ПАЧКА 4: ВАРИАНТЫ 17–20
% ====================================================================

% === ЗАДАЧА 1 ===
% Тег: #МКТ #КИМ_07 #ВАРИАНТ_19
\begin{taskbasic}[code={\label{task:mkt_v19_n7}}]
\textbf{Задача 1.} В результате утечки газа из баллона концентрация его молекул уменьшилась в $3$ раза. Во сколько раз увеличилось давление газа в баллоне, если среднеквадратическая скорость теплового движения его молекул увеличилась в $2$ раза?

\kim{7}
\end{taskbasic}
\answer{1{,}33}

% === ЗАДАЧА 2 ===
% Тег: #МКТ #КИМ_07 #ВАРИАНТ_20
\begin{taskbasic}[code={\label{task:mkt_v20_n7}}]
\textbf{Задача 2.} В сосуде постоянного объёма абсолютную температуру гелия увеличили в $4$ раза, выпустив при этом $\frac{1}{4}$ газа из сосуда. Во сколько раз увеличилось в результате этого давление газа в сосуде?

\kim{7}
\end{taskbasic}
\answer{3}

% ====================================================================
% ФАЙЛ: tasks/02_mkt/mkt_fundamental_equations_bank_part5.tex
% ТЕМА: ОСНОВНОЕ УРАВНЕНИЕ МКТ. АБСОЛЮТНАЯ ТЕМПЕРАТУРА. КИНЕТИЧЕСКАЯ ЭНЕРГИЯ
% ПАЧКА 5: ВАРИАНТЫ 21–30
% ====================================================================

% === ЗАДАЧА 1 ===
% Тег: #МКТ #КИМ_07 #ВАРИАНТ_23
\begin{taskbasic}[code={\label{task:mkt_v23_n7}}]
\textbf{Задача 1.} При уменьшении абсолютной температуры средняя кинетическая энергия поступательного теплового движения молекул разреженного криптона уменьшилась в $1{,}5$ раза. На сколько градусов Цельсия уменьшилась температура газа, если его начальная абсолютная температура составляла $1200$~К?

\kim{7}
\end{taskbasic}
\answer{400}

% === ЗАДАЧА 2 ===
% Тег: #МКТ #КИМ_07 #ВАРИАНТ_24
\begin{taskbasic}[code={\label{task:mkt_v24_n7}}]
\textbf{Задача 2.} При уменьшении абсолютной температуры на $100$~К средняя кинетическая энергия поступательного теплового движения молекул ксенона уменьшилась в $1{,}5$ раза. Какова начальная абсолютная температура газа?

\kim{7}
\end{taskbasic}
\answer{300}